\documentclass[11pt]{article}
\usepackage[margin=2.5cm]{geometry}
\emergencystretch=3em
\usepackage{amsmath,amssymb}
\usepackage{amsthm}
\usepackage{booktabs}
\usepackage{graphicx}
\usepackage{float}
\usepackage{hyperref}
\hypersetup{colorlinks=true, linkcolor=black, citecolor=black, urlcolor=black}

\newtheorem{proposition}{Proposizione}
\newtheorem{definition}{Definizione}
\renewcommand{\proofname}{Dimostrazione}

\title{\Large Trasformazione della Liquidit\`a: perch\'e un ETF pu\`o essere liquido\\anche se detiene asset illiquidi\\[4pt]
\large Sintesi dell'analisi NAV smoothing (LQD vs HYG)}
\author{Tommaso Zeri \\ \small Advanced Fixed Income and Credit --- Analisi supplementare \\ \small Tema: \emph{Market microstructure insights}}
\date{Settembre 2026}

\begin{document}
\maketitle

\begin{abstract}
\noindent Un ETF obbligazionario detiene centinaia di obbligazioni corporate che, individualmente, si scambiano poco e per appuntamento. Eppure l'ETF stesso si scambia in borsa in modo continuo, con spread stretti. Questa nota sintetizza un'analisi complementare al progetto principale (che misura \emph{quando} l'arbitraggio degli Authorized Participants si rompe, tramite il premio/sconto): qui misuriamo \emph{perch\'e} il prezzo dell'ETF pu\`o restare pulito anche quando il suo Net Asset Value (NAV) non lo \`e, usando l'autocorrelazione seriale dei rendimenti come prova statistica diretta della differenza fra i due processi di formazione del prezzo. Validiamo lo stimatore su dati sintetici con un parametro di staleness noto, poi lo applichiamo al NAV storico reale di LQD (investment-grade) e HYG (high-yield), trovando che il NAV di HYG \`e sistematicamente pi\`u stantio di quello di LQD -- coerente con la minore liquidit\`a delle obbligazioni high-yield.
\end{abstract}

% ============================================================================
\section{La domanda e l'idea}

Il progetto principale spiega \emph{quando} il collegamento fra prezzo dell'ETF e NAV si rompe, tramite il premio/sconto $\pi_t = P_t - NAV_t$: quando gli Authorized Participants (AP) non riescono pi\`u ad arbitraggiare, $\pi_t$ si allarga. Questa nota affronta una domanda logicamente precedente: \textbf{perch\'e, in condizioni normali, il prezzo dell'ETF pu\`o essere un segnale di liquidit\`a pulito, anche se le obbligazioni sottostanti non lo sono?}

L'idea: il NAV eredita l'illiquidit\`a delle obbligazioni sottostanti sotto forma di \emph{autocorrelazione seriale positiva} dei suoi rendimenti (rendimenti ``stantii'', che riflettono informazione con ritardo). Il prezzo dell'ETF, formandosi continuamente tramite trading sul secondario, non ha questo meccanismo -- dovrebbe mostrare autocorrelazione vicina a zero. Il divario fra le due \`e la prova, misurabile senza mai vedere l'order book delle obbligazioni, che l'ETF \`e liquido \emph{nonostante}, non \emph{perch\'e}, le sue holding lo siano.

% ============================================================================
\section{Il modello di NAV smoothing}
\label{sec:model}

L'idea, sviluppata originariamente per hedge fund e private equity illiquidi (Getmansky, Lo e Makarov, 2004), si applica identica a un bond ETF. Sia $r^\ast_t$ il rendimento ``vero'' (non osservabile) del paniere di obbligazioni al giorno $t$. Il rendimento \textbf{riportato} nel NAV non \`e $r^\ast_t$, ma una media pesata dei rendimenti veri recenti,
\[
r^{NAV}_t = \sum_{j=0}^{k-1} \theta_j\, r^\ast_{t-j}, \qquad \theta_j \ge 0,\quad \sum_{j=0}^{k-1}\theta_j = 1,
\]
perch\'e quando un'obbligazione non scambia oggi, il pricing desk usa un matrix price basato sull'ultimo prezzo osservato -- informazione vecchia che ``trapela'' nel NAV di oggi.

\paragraph{Perch\'e questo crea autocorrelazione: un esempio minimo.} Il caso pi\`u semplice possibile: met\`a peso su oggi, met\`a su ieri,
\[
r^{NAV}_t = 0.5\, r^\ast_t + 0.5\, r^\ast_{t-1}.
\]
Con rendimenti veri (indipendenti) nei giorni 1--4 pari a $+2\%,\ -1\%,\ +3\%,\ +1\%$:
\begin{align*}
r^{NAV}_2 &= 0.5(-1\%)+0.5(2\%) = 0.5\%, \\
r^{NAV}_3 &= 0.5(3\%)+0.5(-1\%) = 1.0\%, \\
r^{NAV}_4 &= 0.5(1\%)+0.5(3\%) = 2.0\%.
\end{align*}
Si noti che $r^{NAV}_3$ e $r^{NAV}_4$ condividono lo stesso ingrediente, lo shock $+3\%$ del giorno 3 -- una volta come ``oggi'' in $r^{NAV}_3$, di nuovo come ``ieri'' in $r^{NAV}_4$. Quello shock non viene consumato in un colpo solo: met\`a finisce nel NAV di oggi, l'altra met\`a trascina in alto anche il NAV di domani. \`E esattamente questa sovrapposizione di ingredienti fra osservazioni consecutive a generare autocorrelazione positiva. Senza smoothing ($r^{NAV}_t = r^\ast_t$), non c'\`e nessun ingrediente condiviso fra un giorno e il successivo, e l'autocorrelazione \`e nulla.

\paragraph{Caso parametrico generale.} Pesi a decadimento geometrico, $\theta_j = (1-\phi)\phi^{\,j}$, $\phi\in[0,1)$ il \textbf{parametro di smoothing} ($\phi=0$: nessuno smoothing; $\phi\to 1$: smoothing molto persistente).

\begin{proposition}[Autocorrelazione di ordine 1]
Nel limite $k\to\infty$, $r^{NAV}_t$ \`e un processo $MA(\infty)$ con autocovarianza $\gamma(h) = \sigma^2\sum_j\theta_j\theta_{j+h}$. Per i pesi geometrici sopra,
\[
\gamma(h) = \sigma^2\,\frac{(1-\phi)^2}{1-\phi^2}\,\phi^{\,h}
\quad\Longrightarrow\quad
\rho_1 := \frac{\gamma(1)}{\gamma(0)} = \phi.
\]
\end{proposition}

L'autocorrelazione di ordine 1 dei rendimenti del NAV \`e \emph{esattamente} $\phi$: un'identit\`a matematica esatta nel modello, non un'analogia, che rende $\hat\rho_1$ uno stimatore diretto dell'intensit\`a di staleness. Il prezzo dell'ETF, trattato come $r^P_t = r^\ast_t + \eta_t$ con $\eta_t$ rumore idiosincratico i.i.d., \`e a sua volta i.i.d. (somma di due processi i.i.d.\ indipendenti): $\rho_1(r^P)=0$ in popolazione.

\paragraph{Variance ratio (diagnostica complementare).} La variance ratio di Lo e MacKinlay (1988), $VR(q) = \mathrm{Var}(\sum_{i=0}^{q-1}r_{t-i}) / (q\,\mathrm{Var}(r_t))$, vale $1$ sotto i.i.d.\ e cresce con $\phi$ per il processo geometrico sopra (l'autocovarianza decade come $\phi^h$, esattamente come un $AR(1)$).

% ============================================================================
\section{Validazione su dati sintetici}

Simuliamo rendimenti veri i.i.d., li smussiamo con $\phi$ noto per ottenere il NAV sintetico, generiamo un prezzo ETF sintetico come rendimento vero pi\`u rumore idiosincratico. La Tabella~\ref{tab:validation} conferma che l'autocorrelazione campionaria del NAV recupera $\phi$ quasi esattamente, mentre quella del prezzo ETF resta vicina a zero per ogni $\phi$.

\begin{table}[H]
\centering
\caption{Recupero di $\phi$ su campioni sintetici lunghi ($n=20{,}000$).}
\label{tab:validation}
\begin{tabular}{ccc}
\toprule
$\phi$ (vero) & $\hat\rho_1(NAV)$ stimato & $\hat\rho_1(ETF)$ stimato \\
\midrule
$0.0$ & $-0.0054$ & $-0.0047$ \\
$0.2$ & $\phantom{-}0.1957$ & $-0.0062$ \\
$0.4$ & $\phantom{-}0.3876$ & $-0.0132$ \\
$0.6$ & $\phantom{-}0.5986$ & $-0.0031$ \\
$0.8$ & $\phantom{-}0.8052$ & $\phantom{-}0.0181$ \\
\bottomrule
\end{tabular}
\end{table}

Un secondo test usa un percorso sintetico con un regime di smoothing che cambia nel tempo ($\phi$ basso $\to$ alto $\to$ basso $\to$ alto $\to$ basso, in 5 tratti di 500 giorni), per verificare che le stime \emph{rolling} (finestra di 60 giorni) seguano il cambiamento di regime.

\begin{figure}[H]
\centering
\includegraphics[width=0.95\textwidth]{../figures/smoothing_rolling_diagnostics.png}
\caption{Pannello superiore: regime di smoothing noto (ground truth). Pannello centrale: $\hat\rho_1$ rolling del NAV (rosso) segue visibilmente il regime noto, mentre quella del prezzo ETF (blu) resta vicina a zero indipendentemente dal regime. Pannello inferiore: la variance ratio rolling mostra lo stesso pattern.}
\label{fig:rolling}
\end{figure}

\paragraph{HMM sul gap di autocorrelazione.} Un HMM Gaussiano (selezione di $K$ via BIC) applicato al vettore $[\hat\rho_1(NAV),\ \hat\rho_1(ETF),\ \hat\rho_1(NAV)-\hat\rho_1(ETF)]$ classifica correttamente i tratti ad alto $\phi$ (Figura~\ref{fig:hmm}).

\begin{figure}[H]
\centering
\includegraphics[width=0.95\textwidth]{../figures/smoothing_hmm_synthetic.png}
\caption{Regimi HMM decodificati sul gap di autocorrelazione. Le bande rosa si allineano con i tratti a $\phi$ alto del ground truth.}
\label{fig:hmm}
\end{figure}

\textbf{Attenzione all'interpretazione.} Un gap grande (banda rosa) \emph{non} \`e un segnale di allarme: significa che il sottostante \`e stantio \emph{e che, nonostante questo, il prezzo dell'ETF resta pulito} -- la trasformazione di liquidit\`a sta funzionando come dovrebbe. In questa simulazione il rumore del prezzo ETF ha varianza costante per costruzione: non abbiamo mai simulato una rottura genuina del meccanismo. Una rottura vera si manifesterebbe diversamente: non con $\hat\rho_1(NAV)$ che sale, ma con $\hat\rho_1(ETF)$ che comincia anch'esso a salire verso quello del NAV -- il gap si chiuderebbe ``dal lato sbagliato''.

% ============================================================================
\section{Applicazione a dati reali: LQD contro HYG}

Usando lo storico reale del NAV (nessuna rete necessaria: dati gi\`a esportati da iShares), calcoliamo $\hat\rho_1$ e $VR(5)$ lato NAV per entrambi i fondi. L'ipotesi economica: le obbligazioni high-yield sono tipicamente meno liquide di quelle investment-grade, quindi il NAV di HYG dovrebbe mostrare pi\`u smoothing.

\begin{table}[H]
\centering
\caption{Diagnostiche di staleness del NAV, dati reali.}
\label{tab:real}
\begin{tabular}{lccc}
\toprule
Fondo & N giorni & $\hat\rho_1(NAV)$, full-sample & $VR(5)$, full-sample \\
\midrule
LQD (investment-grade) & 6{,}079 & $0.0999$ & $1.2201$ \\
HYG (high-yield) & 4{,}897 & $0.3409$ & $1.9647$ \\
\bottomrule
\end{tabular}
\end{table}

\begin{figure}[H]
\centering
\includegraphics[width=0.95\textwidth]{../figures/smoothing_real_nav_rho1.png}
\caption{Autocorrelazione rolling (60gg) del NAV, LQD (sopra) e HYG (sotto). HYG si mantiene visibilmente pi\`u alto in media; si noti anche il picco di LQD attorno al 2020 (stress COVID-19).}
\label{fig:real}
\end{figure}

Il risultato conferma l'ipotesi: $\hat\rho_1(NAV)$ e $VR(5)$ di HYG sono entrambi pi\`u di tre volte quelli di LQD -- evidenza diretta, basata sui soli dati NAV (senza bisogno del prezzo dell'ETF n\'e di dati di order book), che il paniere high-yield \`e strutturalmente pi\`u stantio/illiquido di quello investment-grade.

Un HMM a 2 stati stimato separatamente per fondo classifica LQD nello stato ``stantio'' per il $21.3\%$ dei giorni contro il $15.6\%$ di HYG -- un risultato che, letto ingenuamente, sembrerebbe contraddire la tabella sopra. Non lo fa: i due HMM sono stimati indipendentemente, quindi lo stato ``stantio'' \`e definito \emph{relativamente allo storico di ciascun fondo}, non su una scala assoluta comune. Il confronto sui livelli medi (Tabella~\ref{tab:real}) resta l'evidenza pi\`u diretta; l'HMM per-fondo misura una cosa diversa -- quanto spesso ciascun fondo attraversa episodi insolitamente stantii rispetto al proprio storico, non quale fondo sia pi\`u stantio in assoluto.

% ============================================================================
\section{Conclusioni e limiti}

Il NAV di un bond ETF eredita l'illiquidit\`a delle obbligazioni sottostanti sotto forma di autocorrelazione seriale positiva; il prezzo di borsa, formandosi tramite trading continuo fra investitori sul secondario, non ha ragione teorica -- n\'e, dove verificabile, evidenza -- di ereditare la stessa propriet\`a. L'ETF pu\`o quindi essere liquido nel senso di avere un prezzo pulito e continuamente aggiornato, anche se il valore che quel prezzo insegue si muove in modo strutturalmente pi\`u lento.

Limiti principali: la finestra rolling (60 giorni) \`e una scelta di design con il consueto trade-off bias/varianza; il modello di smoothing geometrico \`e una parametrizzazione conveniente, non l'unica possibile; il confronto completo NAV-contro-prezzo su dati reali richiede l'OHLCV dell'ETF (rete), non eseguito in questa nota -- la sezione NAV-only resta comunque un test reale e autosufficiente; un $\hat\rho_1$ positivo potrebbe in linea di principio riflettere anche altre fonti di autocorrelazione oltre allo smoothing (non isolate qui).

\end{document}
